\begin{abstract}
Theory of Structural Imagination (TSI)은 추론과 행동을 위한 내부 표현이 저장
형식만이 아니라 보존하는 구조와 지원하는 변환에 의해 명세되어야 한다고
제안한다. 이 논문은 그 제안의 통합된 형식 핵심을 제시한다. State는 finite
presented schema 위의 typed relational functor, 공통 tagged carrier 위의
simplicial complex와 metric, monotone filtration, full-support mass, 그리고
action-conditioned tracked transition으로 구성된다. Structural isomorphism과
topology, geometry, relation, filtration, mass의 다섯 survivor layer 보존을
정의한다. Generator relation을 simplicial pair에 연결하는 compatibility
axiom 아래 relational defect의 정량적 상한을 증명하고, structural equivalence,
composition closure, turnover를 포함한 여섯 quantity의 zero criterion,
hard-label metric--measure triangle inequality, 명시적 equivariance와 Lipschitz
가정 아래 quotient-level rollout bound를 확립한다. 이 결과는 지정된 표현
class에 대한 조건부 수학적 결과이며, 인간 cognition이나 unconstrained neural
network가 이를 구현한다는 것을 확립하지 않는다.
\end{abstract}
